\documentclass[11pt,letterpaper]{article}
% Shared layout. Compile each document from this directory.
\usepackage[margin=0.8in,headheight=14pt,headsep=16pt,footskip=22pt]{geometry}
\usepackage[T1]{fontenc}
\usepackage{lmodern,amsmath,amssymb,booktabs,enumitem,titlesec,fancyhdr,microtype,xcolor,hyperref}
\definecolor{inkblue}{RGB}{25,65,100}
\hypersetup{colorlinks=true,linkcolor=inkblue,urlcolor=inkblue,filecolor=inkblue,pdfauthor={AI-assisted course study notes}}
\setlength{\parindent}{0pt}
\setlength{\parskip}{5pt}
\setlength{\emergencystretch}{1.5em}
\setlength{\abovedisplayskip}{7pt}
\setlength{\belowdisplayskip}{7pt}
\titleformat{\section}{\normalsize\bfseries\color{inkblue}}{}{0pt}{}
\titlespacing*{\section}{0pt}{9pt}{3pt}
\setlist[itemize]{leftmargin=15pt,itemsep=2pt,topsep=3pt,parsep=0pt}
\setlist[enumerate]{leftmargin=17pt,itemsep=3pt,topsep=3pt,parsep=0pt}
\pagestyle{fancy}
\fancyhf{}
\fancyhead[L]{\small ECE 750T\enspace /\enspace Responsible AI}
\fancyhead[R]{\small Fall 2026\enspace /\enspace Study notes}
\fancyfoot[R]{\small\thepage}
\renewcommand{\headrulewidth}{0.3pt}
\newcommand{\lesson}[2]{%
  \hypersetup{pdftitle={ECE 750T: #1 -- #2}}%
  {\LARGE\bfseries #1}\par\vspace{3pt}
  {\large #2}\par\vspace{3pt}
  {\footnotesize AI-assisted reconstruction from course slides, not a lecture transcript.
  Discussion reasoning is inferred; classroom responses were not supplied.}\par
  \vspace{3pt}\hrule\vspace{5pt}}
\newcommand{\sources}[1]{\par\vspace{3pt}{\footnotesize\textbf{Slide trail.} #1
  Page numbers refer to PDF pages, not the slides' continuous printed numbering.}\par}
\newcommand{\Prob}{\mathbb{P}}
\newcommand{\E}{\mathbb{E}}
\newcommand{\Yhat}{\widehat{Y}}
\newcommand{\indep}{\mathrel{\perp\!\!\!\perp}}

\begin{document}
\lesson{Lecture 5}{Memorization, disclosure, and private learning}

\textbf{The lesson.} Removing names, fitting a model, or keeping data on devices
does not establish privacy. A meaningful guarantee specifies the protected
contribution, observer, randomized mechanism, parameters, and all releases.
Differential privacy (DP) controls how much one contribution can change output
distributions; it does not promise that nobody can learn facts about a person.

\section*{From anonymized tables to memorizing models}
Quasi-identifiers can link a nameless release to outside information. With the
specified quasi-identifier columns, $k$-anonymity requires every observed pattern
to occur at least $k$ times. Generalizing ages can create matching groups of size
two. But if both rows have asthma, a known member's diagnosis is disclosed even
though their exact row is ambiguous. Sparse movie ratings likewise support linkage
to an otherwise unnamed profile, as the Netflix example illustrates.

A trained model is another possible release channel: exact reproduction of private
training text can expose an example. Capacity also affects generalization. For
$X\in\mathbb R^{n\times d}$ with $d>n$ and rank $n$,
\[
w_0=X^\top(XX^\top)^{-1}y,\qquad w=w_0+v,\quad Xv=0.
\]
There are $d-n$ unconstrained directions, so interpolation does not determine
behaviour away from training inputs. The spurious-feature example shows how a
model can use a majority shortcut while memorizing minority training exceptions.
Zero training error establishes neither good worst-group generalization nor
privacy; this illustration does not claim that every overparameterized model leaks.

\section*{Define the unit before interpreting the guarantee}
The lecture uses add/remove adjacency: $D\sim D'$ differ in one person's
contribution, with one row per person in the count example. The exam reference
sheet specifies one added/removed row. Row-level protection equals person-level
protection only when the contribution unit actually aligns with the person.
A message, a conversation, and a person's entire history are different units.
Replace-one adjacency is another definition and can change sensitivity.

A mechanism is $(\epsilon,\delta)$-DP if for \emph{every ordered} neighboring pair
and \emph{every} output event $S$,
\[
\Prob[M(D)\in S]\le e^\epsilon\Prob[M(D')\in S]+\delta.
\]
Probabilities are over the mechanism's randomness with the datasets fixed.
Pure DP has $\delta=0$; swapping neighbors supplies the reverse inequality.
Smaller $\epsilon$ forces closer distributions; $\delta$ is additive slack in
the event inequality, not a general probability that a person's data will leak.

\section*{Sensitivity determines how much noise is needed}
For a scalar query, $\Delta f=\sup_{D\sim D'}|f(D)-f(D')|$. A count with at
most one contribution per protected person has sensitivity one. A person's total
messages can be unbounded unless contributions are bounded.
The Laplace mechanism releases $f(D)+Z$ with
$Z\sim\operatorname{Lap}(0,b)$ and $b=\Delta f/\epsilon$.

\newpage
\section*{Why the Laplace mechanism works, and what it costs}
At a released value $z$, the ratio of output densities is bounded using the
triangle inequality:
\[
\frac{p_D(z)}{p_{D'}(z)}
=\exp\!\left(\frac{|z-f(D')|-|z-f(D)|}{b}\right)
\le\exp(\Delta f/b)=e^\epsilon.
\]
Integrating over any event yields pure DP. Checking one plotted event gives a
necessary constraint, not a proof for all events. With continuous outputs, reason
with densities or events; the probability of one exact real number is zero.

For the sensitivity-one count, $\E|Z|=1/\epsilon$ and
$\Prob(|Z|\le\log(20)/\epsilon)=.95$. Stronger privacy therefore increases
error for this fixed query. For general sensitivity, replace $1/\epsilon$ by $b$.
Noise needs calibration to the stated contribution unit; adding arbitrary noise
does not by itself establish DP.

\section*{Account for the whole transcript}
\textbf{Post-processing:} rounding a private count, drawing a chart, or querying
an already released DP model preserves its guarantee if no additional private data
are accessed. A fresh private-data lookup is a new access to account for.

\textbf{Composition:} if each conditional release is
$(\epsilon_t,\delta_t)$-DP, the full transcript is
$(\sum_t\epsilon_t,\sum_t\delta_t)$-DP by basic composition. These are valid
bounds, not claims of optimal tightness. Releasing a count $T$ times with fresh
independent $\operatorname{Lap}(0,1)$ noise gives average-noise SD $\sqrt{2/T}$,
but the basic privacy bound grows to $\epsilon=T$. Reusing one released noisy
count is post-processing; repeatedly generating fresh counts exposes more evidence.

For private training, bound contributions (the slide shows gradient clipping),
add calibrated noise, and account across training steps and all disclosed outputs.
Clipping alone is insufficient. Model selection, diagnostics, or repeated training
that use private data also need appropriate accounting; more advanced accountants
can improve on a naive sum without removing the need to track releases.

\section*{Interpret the promise against the right observer}
Under pure DP, the release multiplies odds between fixed neighboring hypotheses
by at most $e^\epsilon$ (and at least $e^{-\epsilon}$), holding other records and
auxiliary information fixed. Thus $\epsilon=1$ permits a factor 2.72, and
$\epsilon=2$ a factor 7.39. This is incremental evidence, not an absolute posterior
probability or a promise against facts already inferable from other information.

In central DP, a trusted curator sees raw data and protects the release. Local DP
randomizes before the collector receives a message, changing the trust boundary.
Federated learning keeps training data on devices but communicates updates; it
supplies no DP guarantee alone. User-level DP needs bounded user influence,
noise, and accounting across rounds.

The homework-chatbot proposal asks whether $(2,10^{-6})$-DP is acceptable. Ask
what unit it protects, who sees raw conversations, whether the budget includes
retraining, and what utility and alternatives justify the release. A bare claim
that a product ``uses DP'' leaves those decisions unresolved.

\textbf{Midterm emphasis (inference).} The reference sheet makes $k$-anonymity,
adjacency, ordered-event DP checks, composition, and post-processing especially
strong priorities. Also prepare sensitivity/Laplace reasoning, central versus local
trust, memorization, and the limits of federated learning; these were taught even
where their formulas are not supplied.

\sources{\href{../L05.pdf}{L05.pdf}, pp.~2--11 (disclosure/memorization),
12--22 (DP and Laplace), 23--28 (accounting/training), 29--32 (deployment and trust).}
\end{document}
